\documentclass[11pt]{article}
\usepackage[margin=25mm]{geometry}
\usepackage{fontspec}
\setmainfont{Latin Modern Roman}
\usepackage{amsmath,amssymb,mathtools}
\usepackage{unicode-math}
\setmathfont{Latin Modern Math}
\usepackage{xurl,hyperref}
\hypersetup{colorlinks=true,urlcolor=blue}
\setcounter{secnumdepth}{0}
\setlength{\parindent}{0pt}
\setlength{\parskip}{5pt}
\setlength{\emergencystretch}{3em}
\newcommand{\sref}[2]{\hyperlink{source:#1:#2}{[S#2]}}
\newcommand{\cataloguescope}[1]{\paragraph{Recorded target.}#1}
\newcommand{\R}{\mathbb R}
\newcommand{\C}{\mathbb C}
\newcommand{\N}{\mathbb N}
\newcommand{\Z}{\mathbb Z}
\newcommand{\Q}{\mathbb Q}
\newcommand{\F}{\mathbb F}
\newcommand{\eps}{\varepsilon}
\begin{document}
% UnsolvedMath ID 20002548; source catalogue code AIM-PHYSICS-0056
\setcounter{section}{20002547}
\section{A logarithmic-volume identity and zero-angular-momentum obstruction}\label{20002548}
{\small UnsolvedMath ID 20002548 · Mathematical Physics}
\subsection{Definitions and mathematical statement}

\paragraph{Shared notation.}$\mathbb N=\{1,2,\ldots\}$; the standard number fields have their usual meanings. All additional conventions are stated locally.


\paragraph{Definitions and notation.}In the Newtonian $n$-body problem, positions $q_i\in\mathbb R^d$ and masses $m_i>0$ obey $m_i\ddot q_i=\nabla_{q_i}U$, where $U(q)=\sum_{i<j}m_im_j/|q_i-q_j|$ and the gravitational constant is normalized to one. Here $T=\tfrac12\sum_i m_i|\dot q_i|^2$, $E=T-U$, and a collision has $q_i=q_j$. The center of mass can be fixed at zero; moment of inertia is $I=\sum_i m_i|q_i|^2$. For four spatial positions, tetrahedron volume is $V(t)=\tfrac16|\det(q_2-q_1,q_3-q_1,q_4-q_1)|$. Remaining bounded away from zero and infinity means that positive finite constants $v_-,v_+$ bound $V(t)$ throughout the stated time interval.

\paragraph{Catalogue wording retained for traceability.}
What restrictions govern the volume of the tetrahedron in the four-body problem? Can an orbit exist with $0<v_-\le V(t)\le v_+<\infty$ forever?
\paragraph{Scope and source qualification.}Forever must specify one or both directions of time and exclude breakdown by collision. The source does not add zero angular momentum or a chosen energy regime, so such conditions are not inserted. The catalogue title/status was paired with a machine-generated research attempt; that attempt is not peer-reviewed evidence of settlement. The statement covers the original source questions rather than only the title’s subcase.
\subsection{Short English statement}
Determine whether a four-body orbit can maintain a positive tetrahedron volume uniformly bounded above and below forever.
\subsection{Research attempt}
\paragraph{Model, date and investigated scope.}
OpenAI GPT-6 Astra (configured model \texttt{gpt-6-astra}), ultra reasoning; generated 11 October 2026. Derived a logarithmic-volume identity and excluded rigidly rotating full-span central tetrahedra in R3. General forever-bounded positive volume remains unresolved.

\paragraph{Sources and literature check.}
Read the supplied complete definition, catalogue wording, scope qualifications and source list for this canonical ID. The original-source citations are inherited as supplied references; no fresh primary-source or novelty search was performed for this attempt. No current open-status claim is inferred from this calculation.

\paragraph{An exact logarithmic-volume identity.}For a nondegenerate spatial four-body configuration set \(X=(q_2-q_1,q_3-q_1,q_4-q_1)\), an invertible \(3\times3\) matrix. On a collision-free interval with nonzero determinant, \(V=|\det X|/6\). Differentiating the determinant gives
\[\frac{d}{dt}\log V=\operatorname{tr}(X^{-1}\dot X),\qquad
\frac{d^2}{dt^2}\log V=\operatorname{tr}(X^{-1}\ddot X)-\operatorname{tr}[(X^{-1}\dot X)^2].
\]
The second trace need not be nonnegative: for a skew-symmetric matrix \(A\), \(\operatorname{tr}(A^2)=-\|A\|_F^2\). Thus a concavity argument based on dropping this term is invalid without further dynamical restrictions.
\paragraph{A rigid-tetrahedron obstruction.}A three-dimensional regular tetrahedron is a central configuration, but it cannot simply rotate rigidly in \(\mathbb R^3\) as a relative equilibrium. More generally, if a full-span central configuration \(a_i\) satisfies \(\ddot q_i=-\lambda q_i\), \(\lambda>0\), and \(q_i=e^{t\Omega}a_i\) with constant skew \(\Omega\), then spanning forces \(\Omega^2=-\lambda I_3\). Every real three-dimensional skew matrix has a nonzero kernel, so this is impossible. A four-dimensional double rotation is not an admissible spatial tetrahedron example.
\paragraph{Gap.}The identity does not rule out nonrigid solutions with volume bounded above and below forever, since shape distortion and velocity terms can compensate. No zero-angular-momentum assumption is inserted into the source's unrestricted question. The first missing step is a global estimate connecting \(X^{-1}\ddot X\) and its deformation rate strongly enough to force collapse/escape or construct a bounded positive-volume orbit. Both time directions and collision-free lifetime still need specification in such a theorem.

\paragraph{Verification and adversarial audit.}
Checked determinant absolute value, matrix derivative order, skew-square sign, odd-dimensional kernel and exclusion of extra angular-momentum hypotheses. This was an LLM adversarial audit, not independent human verification; no theorem prover was used. Prompts A then B applied to this ID: Checked determinant absolute value, matrix derivative order, skew-square sign, odd-dimensional kernel and exclusion of extra angular-momentum hypotheses. Scope remains limited: Control global deformation and force terms to prove an obstruction or construct a nonrigid positive-volume orbit for the required lifetime.

Coordinator checked both logarithmic-determinant derivatives and the skew-matrix sign obstruction; a full-span three-dimensional rigid central configuration cannot satisfy the stated rotation equation. This additional check was model-assisted, not independent human verification.

\paragraph{Tools, process and contributions.}
Python 3 and local shell used to read supplied files and serialize this manuscript; mathematics derived symbolically in the text, with no numerical experiment or formal proof assistant. Contribution prompt A was followed by prompt B. GPT-6 Astra with ultra reasoning drafted the research and performed its stated self-audit; any additional coordinator contribution is identified in the verification and audit. A GPT-6-based Codex coordinator assembled the manuscript and checked repository integration; its exact variant and reasoning configuration were not exposed. No human mathematical contribution, independent human proof check or formal proof certification is supplied. Complete derivations and the consulted source references are included here; private reasoning transcripts are not published.

\paragraph{Final status and first remaining gap.}
\textbf{UNRESOLVED}. Control global deformation and force terms to prove an obstruction or construct a nonrigid positive-volume orbit for the required lifetime.

\section{Completion Estimate}
\noindent\textbf{COMPLETION: 0\%}
The calculation is an elementary benchmark or conditional reduction; it establishes no claimed new portion of the complete source target.

\subsection{Sources}
\begin{itemize}
\item[\hypertarget{source:20002548:1}{S1}] ULAM.AI and the UnsolvedMath Contributors, UnsolvedMath, record 20002548 (AIM-PHYSICS-0056). Catalogue curation under CC BY 4.0; definitions and English exposition added locally.
\url{https://huggingface.co/datasets/ulamai/UnsolvedMath}.
\item[\hypertarget{source:20002548:2}{S2}] American Institute of Mathematics, Variational Methods in Celestial Mechanics; notes by Kuo-Chang Chen, version 22 June 2003, source item 19.
\url{https://aimath.org/WWN/varcelest/varcelest.pdf}.
\end{itemize}
\label{end:20002548}
\end{document}
